% Created 2024-01-10 mié 16:07
% Intended LaTeX compiler: pdflatex
\documentclass[11pt]{article}
\usepackage[utf8]{inputenc}
\usepackage[T1]{fontenc}
\usepackage{graphicx}
\usepackage{longtable}
\usepackage{wrapfig}
\usepackage{rotating}
\usepackage[normalem]{ulem}
\usepackage{amsmath}
\usepackage{amssymb}
\usepackage{capt-of}
\usepackage{hyperref}
\usepackage{minted}

\usepackage{parskip}
\usepackage[utf8]{inputenc} %% For unicode chars
\usepackage{placeins}
\usepackage[margin=2.50cm]{geometry}
\usepackage[T1]{fontenc}
\usepackage{mathpazo}
\linespread{1.05}
\usepackage[scaled]{helvet}
\usepackage{courier}
\hypersetup{colorlinks=true,linkcolor=blue}
\RequirePackage{fancyvrb}
\DefineVerbatimEnvironment{verbatim}{Verbatim}{fontsize=\small,formatcom = {\color[rgb]{0.5,0,0}}}
\usepackage{mdframed}
\BeforeBeginEnvironment{minted}{\begin{mdframed}}
\AfterEndEnvironment{minted}{\end{mdframed}}
\setcounter{secnumdepth}{3}
\date{}
\title{}
\hypersetup{
 pdfauthor={},
 pdftitle={},
 pdfkeywords={},
 pdfsubject={},
 pdfcreator={Emacs 27.1 (Org mode 9.5.5)}, 
 pdflang={Spanish}}
\begin{document}

\setcounter{tocdepth}{3}
\tableofcontents

\setlength\parindent{10pt}

\section{PHP Functions, more advanced topics}
\label{sec:org122bab2}
Let's delve into the more advanced aspects of PHP functions, including
anonymous functions, variable scope, and recursive functions. These
concepts allow for more dynamic and flexible programming in PHP.

\subsection{Anonymous Functions}
\label{sec:orgf630a5a}
Anonymous functions, are functions without a specified name. They are
most useful as the value of callback parameters, but they have many
other uses.

\begin{itemize}
\item \textbf{Syntax and Usage}
\end{itemize}

Anonymous functions are declared with the \texttt{function} keyword, followed
by a set of parentheses to accept parameters and curly braces to define
the function body.

\begin{minted}[]{php}
<?=
$greet = function($name) {
    echo "Hello, $name!";
};

$greet("World"); // Outputs: Hello, World!
\end{minted}

\begin{itemize}
\item \textbf{Features of Anonymous Functions}
\end{itemize}


\begin{itemize}
\item \textbf{Variable Inheritance}: Anonymous functions have the ability to
inherit variables from the parent scope. This is achieved using the
\texttt{use} keyword.

\begin{minted}[]{php}
<?=
$message = 'Hello';
$example = function() use ($message) {
    echo $message;
};
$example(); // Outputs: Hello
\end{minted}

\item \textbf{Callbacks}: They are often used as callback functions in array
operations, such as \texttt{array\_map} or \texttt{array\_filter}.

\begin{minted}[]{php}
<?=
$numbers = [1, 2, 3, 4, 5];
$squared = array_map(function($n) { return $n * $n; }, $numbers);
print_r($squared); // Outputs: Array ( [0] => 1 [1] => 4 [2] => 9
                   // [3] => 16 [4] => 25 )
\end{minted}
\end{itemize}

\subsubsection{Closures}
\label{sec:org877c9ee}
A closure is an anonymous function that can capture variables local to the
scope it was created in. In PHP, all anonymous functions are closures and
they are implemented using the Closure class. They can specify variables
to be captured with a use clause in the function header.
\begin{minted}[]{php}
<?=
 $x = 1, $y = 2;
 $myClosure = function($z) use ($x, $y)
 {
   return $x + $y + $z;
 };
 echo $myClosure(3); // "6"
\end{minted}

\subsection{Variable Scope}
\label{sec:orge1662b3}
\subsubsection{Definition}
\label{sec:org110a01b}
Variable scope determines the visibility and accessibility of a variable
within a PHP program.
\subsubsection{\emph{Basic} Scopes}
\label{sec:orgb13e4c9}
PHP has three basic scopes:

\begin{itemize}
\item \textbf{Local scope:} Variables declared within a function are local to that
function and not accessible outside of the function. This means that
the \texttt{sum} variable defined in the \texttt{calculateSum()} function cannot be
accessed outside of the function.
\end{itemize}

\begin{minted}[]{php}
<?=
function calculateSum($num1, $num2) {
    $sum = $num1 + $num2;
    echo "The sum is: $sum";
}

$result = calculateSum(10, 20);
echo $result;// This will output an error because 'sum'
             // is not defined outside the function
\end{minted}

\begin{itemize}
\item \textbf{Global scope:} Variables declared outside of any function are global
and can be accessed from anywhere in the program. This means that the
\texttt{age} variable declared outside of any functions can be accessed from
within the \texttt{printAge()} function.
\end{itemize}

\begin{minted}[]{php}
<?php
$age = 30;

function printAge() {
    echo "Your age is: $age";
}

printAge(); // This will print "Your age is: 30"
\end{minted}

\begin{itemize}
\item \textbf{Superglobal variables:} Superglobal variables are a special type of
variable that is available globally, even within functions. They are
used to retrieve data from the request sent by the client (browser) to
the server. Some of the most common superglobal variables include:

\begin{itemize}
\item \texttt{\$\_GET}: Contains data from the \texttt{\$\_GET} superglobal variable, which
is used to pass data to the script through the URL query string.
\item \texttt{\$\_POST}: Contains data from the \texttt{\$\_POST} superglobal variable,
which is used to pass data to the script through the form POST
method.
\item \texttt{\$\_REQUEST}: Contains data from both the \texttt{\$\_GET} and \texttt{\$\_POST}
superglobal variables, as well as additional data from other methods
such as PUT, DELETE, and PATCH.
\item \texttt{\$\_SERVER}: Contains information about the server environment, such
as the client's IP address, the request method, and the server's
protocol.
\item \texttt{\$\_FILES}: Contains information about uploaded files, such as the
file name, file type, and file size.
\item \texttt{\$\_COOKIE}: Contains data from cookies that are sent by the script
to the client's browser and stored on the client's computer.
\end{itemize}
\end{itemize}

\subsubsection{Advanced Scope Concepts}
\label{sec:orgb31221e}
In addition to the basic scopes, there are a few more advanced scope
concepts in PHP:

\begin{itemize}
\item \textbf{Static scope:} Static variables maintain their value between function
calls. They are declared within a function using the \texttt{static} keyword.
\end{itemize}

\begin{minted}[]{php}
<?=
function testStatic() {
    static $z = 0;
    $z++;
    echo $z;
}
testStatic(); // Outputs: 1
testStatic(); // Outputs: 2
\end{minted}

\begin{itemize}
\item \textbf{Closure scope:} Closures are functions that can capture variables
local to the scope they were created in. In PHP, all anonymous
functions are closures and they are implemented using the Closure
class. They can specify variables to be captured with a \texttt{use} clause
in the function header.
\end{itemize}

\begin{minted}[]{php}
<?=
function testClosure() {
    $x = 1, $y = 2;
    $myClosure = function($z) use ($x, $y) {
        return $x + $y + $z;
    };
    echo $myClosure(3); // "6"
}
testClosure();
\end{minted}

\begin{itemize}
\item \textbf{Dynamic scope:} Dynamic scope is a feature of PHP that allows
variables to be accessed from scopes other than their own. This is
typically done with the \texttt{\$GLOBALS} array, but it can also be done with
other techniques such as \texttt{eval()} and \texttt{extract()}. Dynamic scope can
be dangerous because it can lead to unexpected behavior, so it is
generally not recommended.
\end{itemize}

\subsubsection{Conclusion}
\label{sec:org5546b0f}
Understanding variable scope is essential for writing good, maintainable
PHP code. It is important to be aware of the different scopes and how
they interact with each other. By using the right scope for your
variables, you can make your code more readable, reusable, and less
prone to errors.

\subsection{Recursive Functions}
\label{sec:org5a77969}
A recursive function is a function that calls itself during its
execution. This technique is best used for solving problems that can be
broken down into simpler, repetitive tasks.

\subsubsection{Example of a Recursive Function}
\label{sec:org4b081ce}
A classic example is the calculation of the factorial of a number.

\begin{minted}[]{php}
<?=
  function factorial($n) {
      if ($n == 0) {
          return 1;
      } else {
          return $n * factorial($n - 1);
      }
  }

  echo factorial(5); // Outputs: 120
\end{minted}

In this example, \texttt{factorial()} calls itself with a decremented value
until it reaches the base case (\texttt{\$n =} 0=).

\subsubsection{Conclusion}
\label{sec:org3556437}
Understanding these advanced aspects of PHP functions opens up a myriad
of possibilities for writing more efficient and effective code.
Anonymous functions provide flexibility, especially in functional
programming and event-driven programming. Understanding variable scope
is crucial for managing data within and outside functions, while
recursive functions are essential for solving certain types of iterative
problems efficiently.

With these tools in your PHP arsenal, you are well-equipped to tackle
more complex programming challenges and to write code that is both
powerful and elegant.

\section{Functions other advanced topics}
\label{sec:org4442172}
In recent versions of PHP, two significant additions to the language are
arrow functions and generators. These features enhance the functionality
and readability of PHP code, particularly in scenarios involving
callbacks and iterative operations.

\subsection{Arrow Functions}
\label{sec:org804841e}
Introduced in PHP 7.4, arrow functions provide a more concise syntax for
writing anonymous functions. They are similar to anonymous functions
(closures) but with a simpler syntax that's useful for defining simple
callback functions.

\subsubsection{Syntax}
\label{sec:orgaa648ae}
Arrow functions use the \texttt{fn} keyword. The basic syntax is:

\begin{minted}[]{php}
<?=
$arrowFunction = fn($arguments) => $expression;
\end{minted}

\subsubsection{Characteristics}
\label{sec:org2c3fe58}
\begin{itemize}
\item \textbf{Concise Syntax}: Arrow functions allow for a more compact syntax
compared to traditional anonymous functions.
\item \textbf{Automatic Variable Inheritance}: Variables used inside the arrow
function are automatically captured by value from the parent scope.
\end{itemize}

\subsubsection{Example}
\label{sec:org9f7c66b}
Using an arrow function to square numbers in an array:

\begin{minted}[]{php}
<?=
$numbers = [1, 2, 3, 4, 5];
$squared = array_map(fn($n) => $n * $n, $numbers);
// $squared = [1, 4, 9, 16, 25]
\end{minted}

\subsection{Generators}
\label{sec:orgb0f8693}
Generators, introduced in PHP 5.5, are a special type of function that
allow you to iterate through data without needing to create an array in
memory. They are particularly useful for handling large datasets or
streams of data.

\subsubsection{Syntax}
\label{sec:org3723843}
A generator is defined like a regular function, but instead of returning
a value, it yields values using the \texttt{yield} keyword.

\begin{minted}[]{php}
<?=
function myGenerator() {
    yield 'value1';
    yield 'value2';
    // ...
}
\end{minted}

\subsubsection{Characteristics}
\label{sec:org663b641}
\begin{itemize}
\item \textbf{Memory Efficiency}: Instead of loading an entire dataset into memory,
a generator yields one value at a time, making it more
memory-efficient.
\item \textbf{State Preservation}: Generators remember their state between
successive calls.
\end{itemize}

\subsubsection{Example}
\label{sec:orgad15ad9}
Using a generator to iterate over a range of numbers:

\begin{minted}[]{php}
<?=
function generateNumbers($start, $end) {
    for ($i = $start; $i <= $end; $i++) {
        yield $i;
    }
}

foreach (generateNumbers(1, 5) as $number) {
    echo $number . ' ';
}
// Output: 1 2 3 4 5
\end{minted}

\subsection{Conclusion}
\label{sec:org4618ad3}
Arrow functions and generators are powerful additions to PHP, enhancing
the language's capabilities for functional programming and data
iteration. Arrow functions provide a more concise and readable syntax
for small functions or callbacks, while generators offer an efficient
way to handle large datasets or streams, conserving memory and
processing power.
\end{document}
